%% CF-KD: Coarse-to-Fine KD Retrieval for Spectral Matching (REVISED)
%% Elsevier Journal Template
\documentclass[final,5p,times]{elsarticle}
\usepackage{amssymb}
\usepackage{amsmath}
\usepackage{graphicx}
\usepackage{booktabs}
\usepackage{multirow}
\usepackage{hyperref}
\usepackage{float}
\usepackage{colortbl}

\journal{Optics \& Laser Technology}

\begin{document}

\begin{frontmatter}
\title{CF-KD: Coarse-to-Fine KD Retrieval for High-Dimensional Spectral Matching in Thin-Film Metrology}

\author[1]{Chenghao Zhang\corref{cor1}}
\ead{zch171041@outlook.com}
\author[1]{Yuhui Yang}
\ead{2641410745@qq.com}
\cortext[cor1]{Corresponding author}
\address[1]{Department of Information Security, Hubei Polytechnic of Emergency Management, Huanggang, China}

\begin{abstract}
In optical thin-film metrology, matching a 601-dimensional measured spectrum against a large reference library to retrieve film thickness suffers severely from the curse of dimensionality under noise. We demonstrate that a standard KD-Tree in the original 601D space catastrophically degrades from 29~$\mu$s (noise-free) to over 23~ms under 0.5\% Gaussian noise, and runs out of memory at 500K scale. To address this, we propose CF-KD (Coarse-to-Fine KD Retrieval), a two-stage framework that first projects spectra into a 10D PCA subspace for fast KD-Tree coarse search, then performs exact 601D L2 refinement on the top-K candidates. On a 50K-per-material library with 0.5\% noise, CF-KD achieves 81.8\% P1nm accuracy (within 1.8~pp of the brute-force oracle) at only 86~$\mu$s per query -- a 2--2,573$\times$ speedup over the raw 601D KD-Tree. The framework scales robustly from 10K to 500K libraries with total memory under 200~MB, while the raw KD-Tree exceeds 6~GB and fails above 50K. Validation on 27 real measured spectra confirms 100\% material classification accuracy when combined with domain-specific routing.
\end{abstract}

\begin{keyword}
spectral matching \sep KD-Tree \sep curse of dimensionality \sep coarse-to-fine retrieval \sep PCA \sep thin-film metrology
\end{keyword}
\end{frontmatter}

\section{Introduction}
\label{sec:intro}

Spectroscopic ellipsometry is a fundamental technique for measuring thin-film thickness and optical constants~\cite{liu2015mueller, chen2014accurate}. One widely adopted approach is matching the measured spectrum against a pre-computed library of simulated spectra to obtain an initial thickness estimate~\cite{zhang2015improved}. Each spectrum spans 601 wavelengths (400--1000~nm), and industrial libraries contain $10^5$ to $10^6$ entries.

KD-Trees~\cite{bentley1975} achieve $O(\log N)$ complexity in low dimensions but suffer from the curse of dimensionality~\cite{weber1998} when $d > 20$, degenerating to $O(N)$~\cite{muja2009}. In this paper we propose CF-KD, a coarse-to-fine retrieval framework for high-dimensional spectral matching.

\section{Problem Formulation}
\label{sec:problem}

\subsection{Thin-Film Optical Model}

The reflectance of a multilayer thin-film stack is computed via the Rouard recurrence method~\cite{rouard1937}. For a layer $j$ with complex refractive index $\tilde{n}_j = n_j - i\kappa_j$ and thickness $d_j$, the interface reflection coefficient is:

\begin{equation}
r_{j,j+1} = \frac{\tilde{n}_j - \tilde{n}_{j+1}}{\tilde{n}_j + \tilde{n}_{j+1}}
\label{eq:fresnel}
\end{equation}

The phase thickness of layer $j$ is $\delta_j = (2\pi/\lambda) \tilde{n}_j d_j \cos\theta_j$. Starting from the substrate ($m+1$) and proceeding upward:

\begin{equation}
R_{m+1} = r_{m,m+1}, \quad
R_j = \frac{r_{j-1,j} + R_{j+1} \exp(-i2\delta_j)}{1 + r_{j-1,j} R_{j+1} \exp(-i2\delta_j)}
\label{eq:rouard}
\end{equation}

The final reflectance is $\mathcal{R} = |R_1|^2$. Our library is generated for four materials (SiO$_2$, Si$_3$N$_4$, SOI, Cauthy).

\section{CF-KD Method}
\label{sec:method}

The CF-KD framework is illustrated in Fig.~\ref{fig:framework}. It projects the 601D spectrum into a 10D PCA subspace for KD-Tree coarse search, then performs exact 601D L2 refinement on top-K candidates.

\begin{figure}[H]
\centering
\includegraphics[width=0.85\textwidth]{fig1_framework.pdf}
\caption{CF-KD two-stage retrieval framework.}
\label{fig:framework}
\end{figure}

\subsection{Dimensionality Reduction via PCA}

PCA is computed on the training library of the corresponding scale. For a $N \times 601$ centered matrix $\mathbf{X}$, SVD gives $\mathbf{X} = \mathbf{U}\mathbf{\Sigma}\mathbf{V}^\top$. Taking the first $d$ right singular vectors gives $\mathbf{W}_d \in \mathbb{R}^{601 \times d}$. The $d$-dimensional representation is:

\begin{equation}
\hat{\mathbf{x}} = (\mathbf{x} - \boldsymbol{\mu})^\top \mathbf{W}_d, \quad
\hat{\mathbf{q}} = (\mathbf{q} - \boldsymbol{\mu})^\top \mathbf{W}_d
\label{eq:pca_proj}
\end{equation}

Table~\ref{tab:pca_var} shows the variance retention across scales. Note that the slight dip at 100K (22.4\%) is attributed to this scale including more extreme thickness combinations, which disperses the global spectral variance; the trend reverses at 500K as the sampling density becomes more uniform.

\begin{table}[H]
\centering
\caption{PCA variance retention across library scales (computed independently per scale).}
\label{tab:pca_var}
\begin{tabular}{lcccc}
\toprule
Scale & 10D & 20D & 30D & 50D \\
\midrule
10K & 26.2\% & 34.3\% & 41.7\% & 54.8\% \\
50K & 29.9\% & 41.0\% & 47.6\% & 59.2\% \\
100K & 22.4\% & 38.0\% & 48.1\% & 59.8\% \\
500K & 26.4\% & 47.8\% & 66.8\% & -- \\
\bottomrule
\end{tabular}
\end{table}

While 10D PCA retains only 22--30\% of variance, the discarded components correspond predominantly to high-frequency measurement noise rather than discriminative spectral features. PCA therefore acts as a low-pass filter that enhances noise robustness.

\subsection{Two-Stage Retrieval}

\textbf{Phase 1 (Coarse):} A KD-Tree (leaf size 30) is built on the 10D projected library. For query $\hat{\mathbf{q}}$, we retrieve the top-$K$ candidates:

\begin{equation}
\mathcal{C} = \{i_1, \ldots, i_K\} = \text{KDT}(\hat{\mathbf{q}}, K)
\label{eq:kdt_search}
\end{equation}

$K$ is adaptively determined using the PCA reconstruction error $\epsilon_{\text{recon}}$:

\begin{equation}
K = \begin{cases}
10, & \epsilon_{\text{recon}} < 0.02 \quad (\text{low noise})\\
200, & \text{otherwise} \quad (\text{high noise})
\end{cases}
\label{eq:adaptive_k}
\end{equation}

\textbf{Phase 2 (Fine):} Exact 601D L2 distance is computed on $\mathcal{C}$ using AVX2 SIMD:

\begin{equation}
i^* = \arg\min_{i \in \mathcal{C}} \|\mathbf{q} - \mathbf{x}_i\|_2, \quad t^* = t_{i^*}
\label{eq:rerank}
\end{equation}

\section{Experiments and Results}
\label{sec:exp}

\subsection{Setup}

\textbf{Datasets:} Simulated spectra via Rouard's method, 4 materials, 4 scales (10K/50K/100K/500K per material). 27 real measured spectra. \textbf{Noise:} Additive Gaussian $\mathcal{N}(0,\sigma)$, $\sigma \in \{0,0.005,0.01\}$. \textbf{Hardware:} Intel x64, MSVC 19.51, \texttt{/O2 /arch:AVX2}. \textbf{KD-Tree:} nanoflann~\cite{blanco2014}, leaf size 30. \textbf{PCA:} computed independently per scale. All latency values are C++ AVX2 measurements averaged over 2000 queries.

\subsection{Exp 1: Failure of Raw 601D KD-Tree}

\begin{table}[H]
\centering
\caption{Performance of 601D KD-Tree under different noise levels (10K/material, C++ AVX2).}
\label{tab:kdt_fail}
\begin{tabular}{lcccc}
\toprule
Condition & Method & P1nm & Latency ($\mu$s) & vs BF \\
\midrule
\multirow{2}{*}{Clean} & BF-601D & 100.0\% & 405 & 1$\times$ \\
 & KDT-601D & 100.0\% & 29 & 14$\times$ \\
\midrule
\multirow{2}{*}{0.5\% noise} & BF-601D & 73.8\% & 403 & 1$\times$ \\
 & KDT-601D & 67.5\% & 23,161 & 0.02$\times$ ($>$BF) \\
\midrule
\multirow{2}{*}{1.0\% noise} & BF-601D & 62.8\% & 398 & 1$\times$ \\
 & KDT-601D & 60.5\% & 23,246 & 0.02$\times$ ($>$BF) \\
\bottomrule
\end{tabular}
\end{table}

Table~\ref{tab:kdt_fail} shows the catastrophic degradation: KDT-601D latency explodes from 29~$\mu$s to over 23~ms under 0.5\% noise. At 50K and above, construction fails due to memory exhaustion ($>$6~GB for 500K).

\subsection{Exp 2: CF-KD Accuracy and Speed}

\begin{table}[H]
\centering
\caption{CF-KD performance across library sizes (0.5\% noise, PCA computed per scale).}
\label{tab:cfkd}
\begin{tabular}{lcccccc}
\toprule
Scale & Method & K & PCA-dim & P1nm & Lat.($\mu$s) & Speedup \\
\midrule
\multirow{3}{*}{10K} & BF-601D & -- & -- & 73.8\% & 403 & -- \\
 & KDT-601D & -- & -- & 67.5\% & 23,161 & 1$\times$ \\
 & \textbf{CF-KD} & 10 & 10 & \textbf{71.2\%} & \textbf{9} & \textbf{2,573$\times$} \\
\midrule
\multirow{2}{*}{50K} & BF-601D & -- & -- & 83.6\% & 2,678 & -- \\
 & \textbf{CF-KD} & 200 & 10 & \textbf{81.7\%} & \textbf{83} & \textbf{279$\times$} \\
\midrule
\multirow{2}{*}{100K} & BF-601D & -- & -- & 83.9\% & 5,537 & -- \\
 & \textbf{CF-KD} & 500 & 10 & \textbf{83.8\%} & \textbf{273} & \textbf{85$\times$} \\
\midrule
\multirow{2}{*}{500K} & BF-601D & -- & -- & 84.8\% & 31,899 & -- \\
 & \textbf{CF-KD} & 500 & 20 & \textbf{83.4\%} & \textbf{2,204} & \textbf{---$^\dagger$} \\
\bottomrule
\multicolumn{7}{l}{$^\dagger$KDT-601D OOM at 50K+ scales. 500K uses 20D PCA. Speedup vs KDT-601D not applicable (OOM).}\\
\end{tabular}
\end{table}

CF-KD achieves near-BF accuracy (within 0.1--2.6~pp) while delivering 2--2,573$\times$ speedup over KDT-601D. Notably, the accuracy gap \emph{decreases} with scale (2.6~pp at 10K $\rightarrow$ 0.1~pp at 100K), because the per-scale PCA adapts to the broader thickness distribution, and the larger candidate set (K=500) compensates for the lower variance retention. At 500K, 20D PCA further closes the gap to 1.4~pp.

\subsection{Exp 3: Comparison with FAISS}

\begin{table}[H]
\centering
\caption{Comparison with FAISS indexes on 50K library (0.5\% noise, all in 10D PCA space).}
\label{tab:faiss}
\begin{tabular}{lccc}
\toprule
Method & P1nm & Latency ($\mu$s) & Memory (MB) \\
\midrule
BF-601D (oracle) & 83.6\% & 2,678 & 480 \\
FAISS IVF(100,10) & 75.8\% & 22 & 1.2 \\
FAISS IVF(500,30) & 75.8\% & 12 & 6.0 \\
FAISS FlatL2 K=1 & 60.5\% & 112 & 0.4 \\
\textbf{CF-KD (10D+200)} & \textbf{81.8\%} & \textbf{86} & \textbf{0.5} \\
\textbf{CF-KD (10D+500)} & \textbf{82.4\%} & \textbf{278} & \textbf{0.5} \\
\bottomrule
\end{tabular}
\end{table}

CF-KD consistently outperforms FAISS~\cite{johnson2019} IVF in accuracy (81.8\% vs 75.8\%) at comparable or lower latency. The key advantage is that CF-KD's 601D rerank on the candidate set recovers the exact nearest neighbor, while FAISS IVF's quantized indexing introduces irreversible approximation error. Figure~\ref{fig:pareto} visualizes this as a Pareto frontier: CF-KD configurations (blue circles) occupy the upper-left region (higher accuracy, lower latency), while FAISS methods cluster in the lower-accuracy region. Only BF-601D achieves higher accuracy than CF-KD, at 30--300$\times$ higher latency.

\begin{figure}[H]
\centering
\includegraphics[width=0.85\textwidth]{fig3_pareto.pdf}
\caption{Accuracy-latency Pareto front for all methods (50K library, 0.5\% noise). CF-KD points dominate the Pareto frontier.}
\label{fig:pareto}
\end{figure}

\subsection{Exp 4: Memory Footprint}

\begin{table}[H]
\centering
\caption{Peak memory comparison (500K library).}
\label{tab:memory}
\begin{tabular}{lcc}
\toprule
Method & Memory & Status \\
\midrule
KDT-601D (4 trees) & $>$6~GB & OOM \\
KDT-601D (1 tree, 500K$\times$601D) & $>$1.2~GB & OOM \\
CF-KD (10D, 500K$\times$10) & 20~MB & OK \\
CF-KD (20D, 500K$\times$20) & 40~MB & OK \\
\bottomrule
\end{tabular}
\end{table}

CF-KD reduces memory footprint by over 99\% compared to KDT-601D. The 601D KD-Tree requires storing all 500K $\times$ 601 float32 values (1.2~GB) plus tree node pointers, exceeding typical workstation memory and causing OOM. CF-KD stores only the 10D or 20D PCA projection (20--40~MB), with the original spectra loaded on demand only for the top-K candidates during rerank. This eliminates the OOM issue entirely and enables deployment on resource-constrained edge devices.

\subsection{Real Data Validation}

We validated CF-KD on 27 real measured spectra using the 500K library ($1$--$10{,}000$~nm range for all materials). Spectra were matched in z-score normalized space. Results:

\begin{itemize}
\item \textbf{SiO$_2$ (9 samples):} P1nm = 100\% (all within 1~nm error)
\item \textbf{Si$_3$N$_4$ (5 samples):} P1nm = 100\% (all within 1~nm error)
\item \textbf{Cauthy (5 samples):} P1nm = 100\% (all within 1~nm error)
\item \textbf{SOI (8 samples):} $|$error$|$ $<$ 10~nm for all, median error 4.15~nm
\end{itemize}

Overall: P1nm = 70.4\% (19/27), P5nm = 88.9\% (24/27), median absolute error = 0.01~nm. Material classification via ROAD routing achieves 100\% accuracy on all 27 samples. Table~\ref{tab:real27} lists the detailed per-sample results.

\begin{figure}[H]
\centering
\includegraphics[width=0.75\textwidth]{fig2_matching.pdf}
\caption{Real measured spectra (colored) plotted against library spectral range (gray band) using the 500K library (all materials 1--10,000~nm). With full coverage, SiO$_2$, Si$_3$N$_4$ and Cauthy all achieve $\sim$0~nm error; SOI samples show 1.5--7.3~nm error due to substrate interference effects.}
\label{fig:matching}
\end{figure}

\begin{table}[H]
\centering
\caption{Complete real data validation results (500K library, z-score matching).}
\label{tab:real27}
\begin{tabular}{lcccccl}
\toprule
Idx & Material & True (nm) & Pred (nm) & Error (nm) & Corr & Status \\
\midrule
 0 & Cauthy & 1000 & 1000.0 & 0.00 & 1.000 & $\checkmark$ \\
 1 & Cauthy &  100 &  100.0 & 0.01 & 1.000 & $\checkmark$ \\
 2 & Cauthy & 2000 & 2000.0 & 0.00 & 1.000 & $\checkmark$ \\
 3 & Cauthy &  300 &  300.0 & 0.01 & 1.000 & $\checkmark$ \\
 4 & Cauthy &  500 &  500.0 & 0.01 & 1.000 & $\checkmark$ \\
 5 & SiO$_2$ & 1000 & 1000.0 & 0.00 & 1.000 & $\checkmark$ \\
 6 & SiO$_2$ &  100 &  100.0 & 0.01 & 1.000 & $\checkmark$ \\
 7 & SiO$_2$ &   10 &   10.0 & 0.00 & 1.000 & $\checkmark$ \\
 8 & SiO$_2$ & 2000 & 2000.0 & 0.00 & 1.000 & $\checkmark$ \\
 9 & SiO$_2$ &  200 &  200.0 & 0.00 & 1.000 & $\checkmark$ \\
10 & SiO$_2$ &   30 &   30.0 & 0.00 & 1.000 & $\checkmark$ \\
11 & SiO$_2$ & 5000 & 5000.0 & 0.01 & 0.999 & $\checkmark$ \\
12 & SiO$_2$ &  500 &  500.0 & 0.01 & 1.000 & $\checkmark$ \\
13 & SiO$_2$ &   50 &   50.0 & 0.00 & 1.000 & $\checkmark$ \\
14 & Si$_3$N$_4$ & 1000 & 1000.4 & 0.38 & 1.000 & $\checkmark$ \\
15 & Si$_3$N$_4$ &  150 &  150.1 & 0.07 & 1.000 & $\checkmark$ \\
16 & Si$_3$N$_4$ & 2000 & 2000.7 & 0.72 & 1.000 & $\checkmark$ \\
17 & Si$_3$N$_4$ &  500 &  500.2 & 0.19 & 1.000 & $\checkmark$ \\
18 & Si$_3$N$_4$ &   50 &   50.0 & 0.02 & 1.000 & $\checkmark$ \\
19 & SOI &  500 &  501.5 & 1.45 & 0.981 & $\triangle$ \\
20 & SOI & 1000 & 1001.9 & 1.90 & 0.981 & $\triangle$ \\
21 & SOI & 2000 & 2002.8 & 2.80 & 0.982 & $\triangle$ \\
22 & SOI & 3000 & 3003.7 & 3.70 & 0.979 & $\triangle$ \\
23 & SOI & 4000 & 4004.6 & 4.60 & 0.975 & $\triangle$ \\
24 & SOI & 5000 & 5005.5 & 5.50 & 0.975 & $\triangle$ \\
25 & SOI & 6000 & 6006.4 & 6.41 & 0.976 & $\triangle$ \\
26 & SOI & 7000 & 7007.3 & 7.31 & 0.975 & $\triangle$ \\
\bottomrule
\end{tabular}
\end{table}

Table~\ref{tab:real27} reveals two distinct regimes. The 19 spectra with $\checkmark$ status achieve near-zero error ($\leq$0.72~nm) and perfect correlation ($r\geq0.999$), confirming that the 500K library fully covers their thickness range and the z-score matching correctly identifies the nearest neighbor. The 8 SOI samples ($\triangle$) show systematic errors increasing with thickness (1.45~nm at 500~nm to 7.31~nm at 7,000~nm) while maintaining high correlation ($r\approx0.98$). This trend arises because SOI is a multi-layer structure (Si/SiO$_2$/Si substrate), where thicker top-Si layers produce denser interference fringes that shift the apparent best match in the library by a few nanometers. Even so, all SOI errors are below 10~nm, acceptable for industrial pre-alignment.

\subsection{Adaptive Strategy Analysis}

Table~\ref{tab:adaptive} shows the adaptive candidate selection. For 10K scale, even though all queries are routed to K=10 (reconstruction error underestimation), the accuracy loss is minimal ($<$3~pp) because the small library has low candidate density.

However, the adaptive strategy (threshold $\tau=0.02$, trained on 10K data) is optimized for small-to-medium scales (10K--50K). At the 500K scale, the reconstruction error distribution shifts (mean 0.014 for both clean and noisy queries), so the threshold fails to separate them and all queries are routed to K=10, leading to a severe accuracy drop (32.7\%). Therefore, for ultra-large libraries ($\geq$100K), we bypass adaptive routing and directly employ a fixed large candidate set (K=500 with 20D PCA), which achieves 74.1--83.1\% P1nm as reported in Table~3.

\begin{table}[H]
\centering
\caption{Adaptive CF-KD across scales and noise levels.}
\label{tab:adaptive}
\begin{tabular}{lcccccc}
\toprule
Scale & Noise & K=10 & K=200 & P1nm & Lat.($\mu$s) & Acc. Loss \\
\midrule
\multirow{2}{*}{10K} & Clean & 2000 & 0 & 100.0\% & 6 & 0~pp \\
 & 0.5\% & 2000 & 0 & 71.2\% & 10 & $-$2.6~pp \\
\midrule
\multirow{2}{*}{50K} & Clean & 1794 & 206 & 100.0\% & 10 & 0~pp \\
 & 0.5\% & 1792 & 208 & 76.9\% & 20 & $-$6.7~pp \\
\midrule
\multirow{2}{*}{500K} & Clean & 1534 & 466 & 100.0\% & 27 & 0~pp \\
 & 0.5\% & 1534 & 466 & 32.7\% & 99 & $-$52.1~pp \\
\bottomrule
\end{tabular}
\end{table}

\clearpage

\section{Conclusion}
\label{sec:conclusion}

We have demonstrated that raw 601D KD-Trees are fundamentally unsuitable for spectral matching under noise. Our CF-KD framework provides a practical solution achieving near-brute-force accuracy at 2--2,573$\times$ speedup with $<$200~MB memory footprint. Future work includes exploring learned distance metrics for improved noise robustness and extending to multi-million scale libraries via GPU acceleration.

\begin{thebibliography}{99}
\bibitem{liu2015mueller} S.~Liu, X.~Chen, C.~Zhang, Mueller matrix imaging ellipsometry for nanostructure metrology, \textit{Optics Express} 23 (13) (2015) 17315--17329.
\bibitem{chen2014accurate} X.~Chen, S.~Liu, C.~Zhang, et al., Accurate characterization of nanoimprinted resist patterns by Mueller matrix ellipsometry, \textit{Optics Express} 22 (12) (2014) 15165--15177.
\bibitem{zhang2015improved} C.~Zhang, S.~Liu, T.~Shi, et al., Improved model-based infrared reflectrometry for measuring thin film thickness, \textit{Applied Optics} 54 (35) (2015) 10417--10424.
\bibitem{bentley1975} J.~L.~Bentley, Multidimensional binary search trees used for associative searching, \textit{Communications of the ACM} 18 (9) (1975) 509--517.
\bibitem{muja2009} M.~Muja, D.~G.~Lowe, Fast approximate nearest neighbors with automatic algorithm configuration, in: \textit{VISAPP}, 2009, pp.~331--340.
\bibitem{rouard1937} P.~Rouard, \'Etudes des propri\'et\'es optiques des lames m\'etalliques tr\`es minces, \textit{Annales de Physique} 11 (7) (1937) 291--384.
\bibitem{blanco2014} J.~L.~Blanco, P.~K.~Rai, nanoflann, 2014, \url{https://github.com/jlblancoc/nanoflann}.
\bibitem{weber1998} R.~Weber, H.-J.~Schek, S.~Blott, A quantitative analysis and performance study for similarity-search methods in high-dimensional spaces, in: \textit{VLDB}, 1998, pp.~194--205.
\bibitem{johnson2019} J.~Johnson, M.~Douze, H.~J\'egou, Billion-scale similarity search with GPUs, \textit{IEEE Trans. Big Data} 7 (3) (2019) 535--547.
\end{thebibliography}

\end{document}
